% !PRINTED 651
\weberpage{651}

%<<CONT
ein. Man bekommt so durch einfache Rechnung die Gleichung
\weq{11}{x^4 - 2\,(B^2c + b)\,x^2 - 4\,B\,a\,c\,x + B^4c^2 - 2\,B^2bc + c = 0.}
Die Darstellung der $k$ durch die Formeln (9) ist aber noch nicht vollst\"andig befriedigend, weil zwischen den darin vorkommenden rationalen Gr\"ossen $a$, $b$, $c$ noch die Relation (4) stattfindet, und wir suchen eine Darstellung durch unabh\"angige Gr\"ossen.

Einen besonderen Fall m\"ussen wir zun\"achst abmachen, n\"amlich $b = 0$, was $a = i$ zur Folge hat. Dann geben die Formeln (9)
\[
\begin{aligned}
k_0 &= C + B\sqrt{c} + \frac{1+i}{\sqrt{2}}\sqrt[4]{c}\\
k_1 &= C - B\sqrt{c} + \frac{1-i}{\sqrt{2}}\sqrt[4]{c}\\
k_2 &= C + B\sqrt{c} - \frac{1+i}{\sqrt{2}}\sqrt[4]{c}\\
k_3 &= C - B\sqrt{c} - \frac{1-i}{\sqrt{2}}\sqrt[4]{c},
\end{aligned}
\]
und die biquadratische Gleichung (11) wird
\[
x^4 - 2\,B^2c\,x^2 - 4\,i\,B\,c\,x + B^4c^2 + c = 0.
\]
Dieser Fall geh\"ort aber nur dann hierher, wenn $i$ im K\"orper $\mathfrak{R}$ enthalten ist, also niemals bei reellen K\"orpern.

Wenn aber $b$ nicht verschwindet, so f\"uhren wir eine neue rationale Gr\"osse $h$ ein, indem wir
\[
b = h\,(1 + a^2), \quad c = h^2\,(1 + a^2)
\]
setzen, wodurch dann die Relation (4) identisch befriedigt ist, und es geben die Formeln (9), wenn man $Bh$ durch $B$ ersetzt,
\weq{12}{
\begin{aligned}
k_0 &= C + B\sqrt{1+a^2} + \sqrt{h\,\bigl(1+a^2+a\sqrt{1+a^2}\bigr)}\\
k_1 &= C - B\sqrt{1+a^2} + \sqrt{h\,\bigl(1+a^2-a\sqrt{1+a^2}\bigr)}\\
k_2 &= C + B\sqrt{1+a^2} - \sqrt{h\,\bigl(1+a^2+a\sqrt{1+a^2}\bigr)}\\
k_3 &= C - B\sqrt{1+a^2} - \sqrt{h\,\bigl(1+a^2-a\sqrt{1+a^2}\bigr)}
\end{aligned}
}
Dieser Ausdruck f\"ur $k_0$ hat noch den Vorzug, dass er, wie man auch die Vorzeichen der darin vorkommenden Quadratwurzeln bestimmen mag, nur vier verschiedene Werthe darstellt.
